% !TeX root = ../informe.tex
% ===========================================================================
%  4. Arquitectura de la solución
%
%  Sección central del entregable E1. Las dos cosas que el documento de
%  alcance pide con nombre propio están aquí: el diagrama de microfrontends
%  y el de microservicios. El modelo de datos va en la sección 5.
%
%  Los diagramas salen de diagramas/workspace.dsl (Structurizr). Expórtalos
%  a PDF en figuras/ con estos nombres:
%      contexto.pdf · contenedores.pdf · componentes-reservas.pdf · despliegue.pdf
% ===========================================================================

\section{Arquitectura de la solución}
\label{sec:arquitectura}

\instruccion{Esta sección responde a una sola pregunta: por qué el sistema
está partido así y no de otra manera. Describir los componentes es la parte
fácil; lo que se evalúa es la justificación. Apóyate en las vistas para
mostrar la estructura y en el registro de decisiones para explicarla.}

\subsection{Visión general}
\label{sec:vision-general}

El sistema aplica \textbf{microservicios en el backend y microfrontends en el
frontend}, con un \id{edge} nginx delante que unifica el origen y un API
gateway que es el único camino del navegador hacia los servicios. Son once
piezas desplegables: cuatro paquetes de frontend, cinco servicios de backend
—cuatro de dominio más el gateway—, el edge y una instancia de PostgreSQL con
tres bases independientes.

El principio que gobierna el reparto es la \textbf{propiedad de los datos}: cada
servicio es dueño de un dominio y de la base que lo sostiene, y nadie más la
lee. De ahí salen las fronteras, no al revés. \id{ms-reservas} necesita el
horario de una cancha y no lo consulta en \id{canchas\_db}: se lo pregunta a
\id{ms-canchas}. Esa restricción es la que obliga a que cada servicio tenga una
responsabilidad completa en lugar de ser una capa de otro.

Dentro de cada microservicio de dominio la estructura es \textbf{hexagonal}: el
modelo y los casos de uso definen puertos, y los adaptadores —web, JPA,
cliente HTTP— dependen de ellos y nunca al revés. Es lo que permite probar las
ocho reglas de negocio con el puerto de salida sustituido por un doble, sin
levantar Spring ni una base de datos. El gateway es la excepción deliberada: es
infraestructura transversal, sin dominio propio que proteger, y su estructura
se queda en \id{AuthenticationFilter} $\rightarrow$ \id{RouteConfig}
$\rightarrow$ \id{ErrorHandler}.

\subsection{Notación}
\label{sec:notacion}

Los diagramas siguen el modelo C4 \cite{ejemplo-web}: contexto, contenedores
y componentes, más una vista de despliegue. Están generados desde una única
descripción en Structurizr DSL (\id{diagramas/workspace.dsl}), de modo que
las cuatro vistas se mantienen consistentes entre sí por construcción.

\subsection{Vistas arquitectónicas}
\label{sec:vistas}

\instruccion{Una subsección por vista. En cada una: el diagrama, y debajo la
prosa que explica lo que el diagrama no puede decir — por qué esa frontera,
qué pasa si un elemento falla, qué se decidió a propósito. Un diagrama sin
texto no comunica; un texto que solo enumera lo que ya se ve en el diagrama,
tampoco.}

\subsubsection{\texorpdfstring{\textsf{V-01}}{V-01} Contexto del sistema}
\label{v:01}

\begin{figure}[H]
  \centering
  \diagrama{contexto}
  \caption{Vista de contexto: el sistema y sus dos tipos de usuario.}
  \label{fig:contexto}
\end{figure}

El sistema resuelve la reserva de canchas deportivas de un club para dos tipos
de usuario. El \textbf{usuario final} consulta la disponibilidad de una cancha
en una fecha, reserva un bloque libre, ve sus reservas y las cancela mientras no
hayan empezado. El \textbf{administrador} gestiona el catálogo de canchas y sus
horarios, cancela cualquier reserva y consulta los indicadores de ocupación. No
hay más actores ni integraciones externas: no existe pasarela de pagos ni envío
de notificaciones, que están fuera del alcance.

\subsubsection{\texorpdfstring{\textsf{V-02}}{V-02} Contenedores}
\label{v:02}

Es la vista principal del documento: contiene a la vez la composición de
microfrontends y la de microservicios.

\begin{figure}[p]
  \thisfloatpagestyle{empty}
  \centering
  \diagramapagina{contenedores}
  \caption{Vista de contenedores: microfrontends, gateway, microservicios y
           bases de datos.}
  \label{fig:contenedores}
\end{figure}

\paragraph*{Frontend.}
Cuatro paquetes npm independientes, integrados en tiempo de ejecución con Module
Federation. El \id{shell} es el host: enruta, mantiene la sesión y carga cada
remote de forma perezosa, solo cuando el usuario entra en su ruta. Cada remote
expone únicamente \id{./App} y se construye y despliega por separado, sin
depender del ciclo de vida de los demás.

Lo único que cruza la frontera entre paquetes es la prop \id{sesion} que baja el
shell y los chunks compartidos de React. No hay imports entre remotes: si los
hubiera, dejarían de ser desplegables por separado y el estilo se quedaría en el
nombre.

\begin{table}[H]
  \centering
  \caption{Composición de microfrontends.}
  \label{tab:microfrontends}
  \begin{tabular}{L{3.6cm} L{6.4cm} L{3.4cm}}
    \toprule
    \textbf{Microfrontend} & \textbf{Responsabilidad} & \textbf{Tipo} \\
    \midrule
    \id{shell}             & Enrutado, sesión, login y registro, y carga de los remotes. & Host   \\
    \id{mf-reservas}       & Disponibilidad, nueva reserva y «mis reservas» del usuario final. & Remote \\
    \id{mf-administracion} & Gestión de canchas, de reservas y de usuarios. & Remote \\
    \id{mf-reportes}       & Los cuatro indicadores de ocupación y demanda. & Remote \\
    \bottomrule
  \end{tabular}
\end{table}

\paragraph*{Backend.}
Cuatro microservicios de dominio y un gateway. Tres de los cuatro tienen base
propia; \id{ms-reportes} \textbf{no tiene ninguna}, y es deliberado: un reporte
de ocupación necesita cruzar reservas con canchas, y la forma cómoda de hacerlo
sería leer las dos bases. Eso rompería la propiedad de los datos por
conveniencia de una consulta. En su lugar agrega vía REST, llamando a
\id{ms-reservas} y \id{ms-canchas} por puertos de salida. Ni siquiera tiene el
driver JDBC en su \id{pom.xml}, de modo que la tentación no está disponible.

\begin{table}[H]
  \centering
  \caption{Composición de microservicios.}
  \label{tab:microservicios}
  \begin{tabular}{L{3.2cm} L{6.4cm} L{3.6cm}}
    \toprule
    \textbf{Microservicio} & \textbf{Responsabilidad} & \textbf{Base de datos} \\
    \midrule
    \id{ms-usuarios} & Registro, autenticación y gestión de usuarios y roles. & \id{usuarios\_db} \\
    \id{ms-canchas}  & Catálogo de canchas, horarios de atención y bloqueos de mantenimiento. & \id{canchas\_db}  \\
    \id{ms-reservas} & Reservas y disponibilidad. Dueño de \rn{01}--\rn{06} y \rn{08}. & \id{reservas\_db} \\
    \id{ms-reportes} & Indicadores de ocupación y demanda; agrega vía REST. & ninguna           \\
    \bottomrule
  \end{tabular}
\end{table}

\paragraph*{Comunicación.}
Todo es HTTP con JSON, síncrono. El navegador solo conoce \id{http://localhost}:
el edge sirve el shell en la raíz, cada remote bajo su prefijo, y reenvía
\id{/api/} al gateway. El edge es el \textbf{único contenedor conectado a las dos
redes} de Docker, y eso es lo que impide que el navegador alcance un
microservicio saltándose el gateway.

Entre servicios hay dos llamadas, y las dos existen para no leer una base ajena:

\begin{itemize}
  \item \id{gateway} $\rightarrow$ \id{ms-usuarios}, para verificar la credencial
        de cada petición.
  \item \id{ms-reservas} $\rightarrow$ \id{ms-canchas}, para comprobar que la
        cancha existe, está activa y admite el bloque antes de persistir.
  \item \id{ms-reportes} $\rightarrow$ \id{ms-reservas} y \id{ms-canchas}, para
        agregar los indicadores.
\end{itemize}

Si un destino no responde, el fallo se traduce a \textbf{503} y no a un 500
genérico: la diferencia es la que permite que la pantalla diga «inténtalo de
nuevo» en lugar de «algo se rompió». En \id{ms-reservas} eso significa que una
reserva no se crea si no se pudo validar la cancha; persistirla y reconciliar
después dejaría reservas sobre canchas inexistentes.

\subsubsection{\texorpdfstring{\textsf{V-03}}{V-03} Componentes de \id{ms-reservas}}
\label{v:03}

Se detalla este servicio y no otro porque es donde viven las reglas de
negocio del sistema (\cref{sec:reglas-negocio}).

\begin{figure}[p]
  \thisfloatpagestyle{empty}
  \centering
  \diagramapagina{componentes-reservas}
  \caption{Vista de componentes de \id{ms-reservas}.}
  \label{fig:componentes-reservas}
\end{figure}

Una petición \id{POST /api/reservas} recorre el hexágono en este orden, y cada
salto tiene un motivo:

\begin{enumerate}
  \item \id{BookingController} (\emph{adaptador de entrada}) traduce HTTP a DTO
        y lee la identidad de \id{X-User-Id} y \id{X-User-Role}. No decide nada:
        no hay en él un solo \id{if} de negocio.
  \item \id{BookingUseCase} (\emph{puerto de entrada}) es lo único que el
        controlador conoce. No sabe qué implementación hay detrás.
  \item \id{BookingService} (\emph{aplicación}) aplica las reglas en orden:
        \textsf{FR-018} —un administrador no reserva—, después la cancha vía
        \id{CourtClientPort}, después \rn{06} con el tope leído por
        \id{ConfigurationRepositoryPort}, y por último \rn{02}. El orden importa:
        si el usuario ya llegó al tope, el bloque concreto da igual y el mensaje
        es más útil.
  \item \id{BookingRepositoryPort} (\emph{puerto de salida}) recibe un
        \id{Booking} de dominio. El servicio no sabe que detrás hay JPA.
  \item \id{BookingRepositoryAdapter} (\emph{adaptador de salida}) traduce a la
        entidad JPA, persiste, y \textbf{captura la violación de la restricción
        \id{EXCLUDE} traduciéndola a \id{SlotAlreadyBookedException}}. Es el
        punto donde una excepción de PostgreSQL deja de existir para el dominio.
\end{enumerate}

El reloj se inyecta en \id{BookingService} en lugar de leerse con
\id{LocalDateTime.now()} allí donde haga falta. Parece un detalle y no lo es:
\rn{04} y la derivación de \id{FINALIZADA} dependen del instante actual, y sin
poder fijarlo una prueba tendría que dormir o depender de la hora a la que se
ejecute.

\subsubsection{\texorpdfstring{\textsf{V-04}}{V-04} Despliegue}
\label{v:04}

\begin{figure}[p]
  \thisfloatpagestyle{empty}
  \centering
  \diagramapagina{despliegue}
  \caption{Vista de despliegue: contenedores Docker, redes y volumen.}
  \label{fig:despliegue}
\end{figure}

Dos redes bridge. La red \id{frontend} contiene el shell y los tres remotes; la
red \id{backend}, el gateway, los cuatro microservicios y PostgreSQL. El
\id{edge} es el único contenedor con un pie en cada una, y por tanto el único
camino entre ellas. Un contenedor de frontend no puede alcanzar
\id{ms-reservas} aunque conozca su nombre: no comparten red.

Los datos persisten en un volumen \id{pgdata}. Los scripts de
\id{infra/postgres/init/} —que crean las tres bases, sus roles sin permisos
cruzados y la extensión \id{btree\_gist}— se ejecutan una sola vez, cuando el
volumen está vacío; por eso \id{docker compose down -v} es lo que fuerza un
arranque limpio.

El orden de arranque se apoya en \id{depends\_on} con
\id{condition: service\_healthy}, sobre el \id{health} de Actuator de cada
servicio y \id{pg\_isready} para PostgreSQL. Sin eso el arranque es una carrera:
Flyway intentaría migrar contra una base que todavía no acepta conexiones.

\subsection{Seguridad y control de acceso}
\label{sec:seguridad}

La autenticación es \textbf{HTTP Basic verificada en cada petición por el
gateway}. No se emite ningún token: no hay JWT, ni sesión de servidor, ni nada
que renovar o revocar. El alcance del proyecto pide autenticación básica con
roles, y añadir OAuth2 sería construir un mecanismo que ningún criterio evalúa.

El recorrido de una petición autenticada es:

\begin{enumerate}
  \item El navegador manda \id{Authorization: Basic} en toda petición a
        \id{/api/**}, salvo las dos rutas públicas —\id{POST /api/auth/login} y
        \id{POST /api/auth/registro}—, que existen precisamente para obtener la
        credencial.
  \item El \id{AuthenticationFilter} del gateway \textbf{descarta primero
        cualquier cabecera \id{X-User-*} que venga del cliente}, verifica la
        credencial contra \id{ms-usuarios} y, si es válida y la cuenta está
        activa, añade \id{X-User-Id} y \id{X-User-Role}.
  \item El microservicio destino lee esas dos cabeceras y aplica \rn{03} y
        \rn{07} sobre ellas. No recibe nunca las credenciales del usuario y no
        vuelve a autenticar.
\end{enumerate}

El orden del segundo punto es lo que hace segura la confianza entre el gateway y
los microservicios. Un microservicio cree lo que dice \id{X-User-Id}
precisamente porque nadie más puede escribirla: si el borrado no fuera lo
primero que ocurre, cualquiera se declararía administrador mandando una
cabecera. La prueba \id{AuthenticationFilterTest} fija ese comportamiento, y el
escenario V5 del \id{quickstart} lo comprueba de extremo a extremo con
\id{curl}.

Las contraseñas se guardan hasheadas con BCrypt y nunca en claro. \id{ms-usuarios}
depende de \id{spring-security-crypto} —la librería de criptografía— y no de
\id{spring-boot-starter-security}: no hay cadena de filtros de seguridad en
ningún microservicio, porque quien autentica es el gateway.

\subsection{Registro de decisiones de arquitectura}
\label{sec:decisiones}

\instruccion{Cada decisión relevante con su alternativa descartada. Una
decisión sin alternativa no es una decisión: es una descripción. Este
registro es lo que distingue un documento de arquitectura de un manual de
usuario del propio sistema. Referencia las decisiones desde el texto con
\texttt{\textbackslash adrref\{01\}}.}

\decision{01}{Restricción de exclusión en PostgreSQL para el solapamiento}
  {\rn{02} tiene que sostenerse también cuando dos personas reservan el mismo
   bloque simultáneamente. Una comprobación en el servicio deja una ventana
   entre la lectura y la escritura por la que caben dos peticiones a la vez.}
  {Se mantiene la comprobación previa en \id{BookingService}, por el mensaje
   claro, y se añade una restricción \id{EXCLUDE USING gist} parcial sobre las
   reservas confirmadas. \id{BookingRepositoryAdapter} traduce su violación a
   la misma excepción de dominio que la comprobación previa.}
  {Validar solo en Java, con un bloqueo pesimista sobre la cancha o una
   transacción \id{SERIALIZABLE}. Se descartó: lo primero serializa todas las
   reservas de una cancha aunque sean de días distintos, y lo segundo obliga a
   una política de reintentos en toda la aplicación para un caso que el motor
   ya sabe resolver en el índice.}
  {Se gana una garantía que no depende de que el código la recuerde, y un 409
   idéntico por los dos caminos. Se paga con una dependencia de PostgreSQL y su
   extensión \id{btree\_gist}: el esquema deja de ser portable a otro motor.}

\decision{02}{Una base por microservicio, con un rol sin permisos cruzados}
  {La independencia de datos entre microservicios se evalúa, y una convención
   que solo vive en la cabeza del equipo se rompe en cuanto alguien tiene prisa.}
  {Tres bases en una única instancia de PostgreSQL, cada una con un rol
   propietario al que se le concede \id{CONNECT} solo sobre la suya, revocándolo
   a \id{PUBLIC}. \id{ms-reportes} no recibe rol ni driver JDBC.}
  {Un esquema compartido con prefijos de tabla, o tres instancias separadas de
   PostgreSQL. Lo primero no impide el acceso cruzado, solo lo desaconseja; lo
   segundo triplica memoria y tiempo de arranque sin ganar nada verificable.}
  {Se gana que una violación del principio falle al conectar en lugar de pasar
   desapercibida en revisión. Se paga con una consulta imposible: cruzar
   reservas y canchas exige una llamada HTTP, no un \id{JOIN}.}

\decision{03}{Un edge que unifica el origen, en lugar de CORS por servicio}
  {Los microfrontends federados piden sus chunks en tiempo de ejecución. Con
   cada remote en su propio puerto, el navegador los ve como orígenes distintos
   y bloquea la carga.}
  {Un nginx delante que sirve el shell en la raíz, cada remote bajo su prefijo y
   reenvía \id{/api/} al gateway. Un solo origen para todo.}
  {Configurar cabeceras CORS permisivas en cada remote y en el gateway. Se
   descartó porque hay que repetirlo en cinco sitios y cada uno es una
   oportunidad de abrir de más.}
  {Se gana que el problema desaparezca en lugar de parchearse, y un único punto
   donde publicar el sistema. Se paga con una pieza más en el despliegue.}

\instruccion{Candidatas a decisión en este proyecto, por si ayudan a
arrancar: una base de datos por microservicio frente a un esquema
compartido; la restricción de exclusión en PostgreSQL frente a validar el
solapamiento en Java; el gateway como punto único frente a que cada
microfrontend llame directamente a los servicios; \id{ms-reportes} agregando
por REST frente a leer las tablas de otros servicios; el edge que unifica el
origen frente a configurar CORS en cada remote.}
